\section*{\texorpdfstring{\S\CurrentExerciseGroup}{Section \CurrentExerciseGroup}}
\phantomsection
\addcontentsline{toc}{section}{\texorpdfstring{Exercises for \S\CurrentExerciseGroup}{Exercises for Section \CurrentExerciseGroup}}

\begin{enumerate}
\def\labelenumi{\arabic{enumi})}
\item
  Let \(\mathrm{T}, \mathrm{T}'\) be two locally compact spaces, \(\mu\) a positive measure on \(\mathrm{T}\) , \(\mu'\) a positive measure on \(\mathrm{T}'\) . Show that if the measure \(\mu'\) is bounded, then the projection \(\mathrm{pr}_1\) of \(\mathrm{T} \times \mathrm{T}'\) onto \(\mathrm{T}\) is a \((\mu \otimes \mu')\) -proper mapping and \(\mathrm{pr}_1(\mu \otimes \mu') = a \cdot \mu\) , where \(a = \mu'(\mathrm{T}')\) . Deduce from this an example of a \((\mu \otimes \mu')\) -negligible subset of \(\mathrm{T} \times \mathrm{T}'\) whose projection on \(\mathrm{T}\) is not \(\mu\) -measurable.
\item
  Let T be the interval \([0,1]\) of R, equipped with the topology induced by that of R, and \(T'\) the interval \([0,1]\) of R, equipped with the discrete topology; let \(\mu\) be Lebesgue measure on T, \(\mu'\) the discrete measure on \(T'\) defined by the mass +1 at every point of \(T'\) , and let \(\nu = \mu \otimes \mu'\) be the product measure on \(X = T \times T'\) .
\end{enumerate}

\begin{enumerate}
\def\labelenumi{\alph{enumi})}
\item
  For every \(t \in \mathrm{T}\) , the set \(\{t\}\) is \(\mu\) -negligible, but the set \(\{t\} \times \mathrm{T}'\) is not \(\nu\) -negligible.
\item
  The `diagonal' \(\Delta\) of \(\mathrm{T} \times \mathrm{T}'\) (the set of points \((t, t)\) , where \(t\) runs over \([0, 1]\) ) is a closed set in \(\mathrm{T} \times \mathrm{T}'\) , locally \(\nu\) -negligible but not \(\nu\) -negligible, such that \(\mu'(\Delta(t)) = 1\) for all \(t \in \mathrm{T}\) and \(\mu(\Delta(t')) = 0\) for all \(t' \in \mathrm{T}'\) (cf.~§3, Exer. 3).
\item
  On the space \(\mathrm{X} \times \mathrm{X} = \mathrm{Y}\) , consider the product measure \(\rho = \nu \times \nu\) . Show that there exists in \(\mathrm{Y}\) a set \(\mathrm{A}\) , locally \(\rho\) -negligible, such that for every \(x \in \mathrm{X}\) , \(\mathrm{A}(x)\) and \(\mathrm{A}^{-1}(x)\) are \(\nu\) -measurable and not locally negligible. (If \(x = (t, t')\) , where \(t \in \mathrm{T}\) , \(t' \in \mathrm{T}'\) , take A to be such that A(x) contains all of the points \((s, t) \in X\) , where \(0 \leqslant s \leqslant 1\) .)
\end{enumerate}

\begin{enumerate}
\def\labelenumi{\arabic{enumi})}
\setcounter{enumi}{2}
\tightlist
\item
  Let T be a locally compact space, \(\mu\) a positive measure on T, f a numerical function \(\geqslant 0\) defined on T. In the product space \(T \times R\) , denote by \(D_{f}\) the set of points \((t, x)\) such that \(0 \leqslant x \leqslant f(t)\) ; on the other hand let \(\nu\) be Lebesgue measure on R.
\end{enumerate}

\begin{enumerate}
\def\labelenumi{\alph{enumi})}
\item
  Show that for \(f\) to be \(\mu\) -measurable, it is necessary and sufficient that \(D_f\) be a measurable set for the product measure \(\lambda = \mu \otimes \nu\) . (To see that the condition is necessary, prove that if \(f\) is \(\mu\) -measurable, then for every compact subset \(K\) of \(T\) , \(D_f \cap (K \times R_+)\) is the union of a \(\lambda\) -negligible set and a countable family of sets of the form \(A \times I\) , where \(A\) is \(\mu\) -measurable and \(I\) is an interval of \(R_+\) . To show that the condition is sufficient, prove that if it is satisfied then, for every compact subset \(K\) of \(T\) , there exists a dense set \(H \subset R_+\) such that, for every \(\alpha \in H\) , \(K \cap f^{-1}([\alpha, +\infty])\) is \(\mu\) -measurable, using Cor. 2 of Prop. 7.)
\item
  Show that, for \(f\) to be \(\mu\) -integrable, it is necessary and sufficient that \(\mathrm{D}_f\) be \(\lambda\) -integrable, and one then has \(\lambda(\mathrm{D}_f) = \int f d\mu\) . Moreover, denoting by \(g\) the decreasing numerical function on \(\mathbf{R}_+\) defined by \(g(t) = \mu\left( f([t, +\infty]) \right)\) (possibly equal to \(+\infty\) for \(t = 0\) (Ch. IV, §5, Exer. 29)), one has \(\int f d\mu = \int_0^{+\infty} g(t) dt\) .
\end{enumerate}

\begin{enumerate}
\def\labelenumi{\arabic{enumi})}
\setcounter{enumi}{3}
\tightlist
\item
  \begin{enumerate}
  \def\labelenumii{\alph{enumii})}
  \tightlist
  \item
    Let \(T, T'\) be two locally compact spaces countable at infinity, \(\mu\) a positive measure on T, \(\mu'\) a positive measure on \(T'\) . Show that for every \((\mu \otimes \mu')\) -measurable mapping f of \(T \times T'\) into R, the function \(F : t \mapsto \int^{*} |f(t, t')| d\mu'(t')\) is \(\mu\) -measurable (consider \(T'\) as the union of an increasing sequence of compact sets).
  \end{enumerate}
\end{enumerate}

\begin{enumerate}
\def\labelenumi{\alph{enumi})}
\setcounter{enumi}{1}
\tightlist
\item
  Suppose, moreover, that for almost every \(t \in T\) (for \(\mu\) ) the function \(f(t, \cdot)\) is \(\mu'\) -integrable and that for almost every \(t' \in T'\) (for \(\mu'\) ) the function \(f(\cdot, t')\) is \(\mu\) -integrable; finally, assume that the function (defined almost everywhere for \(\mu\) )
\end{enumerate}

\(t \mapsto \int f(t, t') d\mu'(t')\) is \(\mu\) -integrable. Show that the function (defined almost everywhere for \(\mu'\) ) \(t' \mapsto \int f(t, t') d\mu(t)\) is then \(\mu'\) -integrable and

\[
\int d \mu^ {\prime} (t ^ {\prime}) \int f (t, t ^ {\prime}) d \mu (t) = \int d \mu (t) \int f (t, t ^ {\prime}) d \mu^ {\prime} (t ^ {\prime}).
\]

(By virtue of \(a\) ), there is a partition of \(\mathbf{T}\) formed by a \(\mu\) -negligible set \(\mathbf{N}\) and a sequence \((\mathrm{K}_n)\) of compact sets such that each of the functions \(f\varphi_{\mathrm{K}_n\times \mathrm{T}'}\) is \((\mu \otimes \mu')\) -integrable. Set

\[
g _ {n} (t ^ {\prime}) = \sum_ {j = 1} ^ {n} \int_ {K _ {j}} f (t, t ^ {\prime}) d \mu (t)
\]

and apply Exer. 20 of §5.)

\begin{enumerate}
\def\labelenumi{\alph{enumi})}
\setcounter{enumi}{2}
\tightlist
\item
  Let \(\mathbf{T}\) be the interval \([0,1]\) of \(\mathbf{R}\) , \(\mu\) Lebesgue measure on \(\mathbf{T}\) . For every integer \(n > 0\) , let \(\mathrm{A}_n^\prime = \left[\frac{1}{2^n}, \frac{3}{2^{n+1}}\right]\) , \(\mathrm{A}_n^{\prime\prime} = \left[\frac{3}{2^{n+1}}, \frac{1}{2^{n-1}}\right]\) , and, in the product space \(\mathbf{T} \times \mathbf{T}\) , let \(\mathrm{B}_n^\prime = \mathrm{A}_n^\prime \times \mathrm{A}_n^\prime\) , \(\mathrm{B}_n^{\prime\prime} = \mathrm{A}_n^{\prime\prime} \times \mathrm{A}_n^{\prime\prime}\) , \(\mathrm{C}_n^\prime = \mathrm{A}_n^\prime \times \mathrm{A}_n^{\prime\prime}\) , \(\mathrm{C}_n^{\prime\prime} = \mathrm{A}_n^{\prime\prime} \times \mathrm{A}_n^\prime\) . Set \(f(t,t') = 4^{n+1}\) in \(\mathrm{B}_n^\prime\) and \(\mathrm{B}_n^{\prime\prime}\) , \(f(t,t') = -4^{n+1}\) in \(\mathrm{C}_n^\prime\) and \(\mathrm{C}_n^{\prime\prime}\) , for every integer \(n > 0\) , and \(f(t,t') = 0\) at the other points of \(\mathbf{T} \times \mathbf{T}\) . Show that the two integrals \(\int d\mu(t') \int f(t,t') d\mu(t)\) and \(\int d\mu(t) \int f(t,t') d\mu(t')\) are defined and equal, but the function \(f\) is not integrable for the measure \(\mu \otimes \mu\) .
\end{enumerate}

¶ 5) Let T, T' be two locally compact spaces, μ a positive measure on T, μ' a positive measure on T'; suppose that every compact subset of T is metrizable. Let f be a mapping of T × T' into a metrizable space G, such that: 1° for every t ∈ T, the mapping t' → f(t, t') is μ'-measurable; 2° for every t' ∈ T', the mapping t → f(t, t') is continuous. Show that, under these conditions, f is (μ ⊗ μ')-measurable. (Observe (using Egoroff's theorem) that for every compact subset K of T, the restriction of f to K × T' is the limit of a sequence of (μ ⊗ μ')-measurable functions.)

¶ 6) a) Let T be a locally compact space, ν a positive measure on T, I an interval of R, μ a positive measure on I. Let A be a subset of the product I × T, such that: 1° for every x ∈ I, the section A(x) ⊂ T of A at x is ν-measurable; 2° the relation x ≤ y implies A(x) ⊂ A(y). Show that A is (μ ⊗ ν)-measurable. (Reduce to the case that I and T are compact and μ is diffuse; consider an increasing sequence (xi)0≤i≤n of points of I, x0 and xn being the end-points of I, such that μ({[}xi, xi+1{]}) ≤ ε for 0 ≤ i ≤ n-1, and form two subsets B,C of I × T, measurable for λ = μ ⊗ ν and such that B ⊂ A ⊂ C and λ(C - B) ≤ εν(t).)

\begin{enumerate}
\def\labelenumi{\alph{enumi})}
\setcounter{enumi}{1}
\item
  Deduce from a) that if \(f\) is a numerical function defined on \(\mathrm{I} \times \mathrm{T}\) , such that: \(1^{\circ}\) for each \(x \in \mathrm{I}\) , \(t \mapsto f(x, t)\) is \(\nu\) -measurable, and \(2^{\circ}\) for each \(t \in \mathrm{T}\) , \(x \mapsto f(x, t)\) is increasing, then \(f\) is \((\mu \otimes \nu)\) -measurable.
\item
  Let \(g\) be a numerical function defined on \(I \times T\) , such that: \(1^{\circ}\) for each \(x \in I\) , \(t \mapsto g(x, t)\) is \(\nu\) -measurable; \(2^{\circ}\) for each \(t \in T\) , \(x \mapsto g(x, t)\) is monotone. Deduce from \(b\) ) that \(g\) is \((\mu \otimes \nu)\) -measurable. (Let \(T_1 \subset T\) be the set of \(t \in T\) such that the mapping \(x \mapsto g(x, t)\) is increasing; show that \(T_1\) is \(\nu\) -measurable. To that end, for every pair \((r_1, r_2)\) of rational numbers belonging to \(I\) such that \(r_1 \leqslant r_2\) , consider the set \(D_{r_1, r_2}\) of points \(t \in T\) such that \(g(r_1, t) \leqslant g(r_2, t)\) and express \(T_1\) as the intersection of sets \(D_{r_1, r_2}\) .)
\end{enumerate}

¶ 7) a) Let T, T' be two locally compact spaces; assume that T' admits a countable base. Let μ be a positive measure on T, μ' a positive measure on T'. Let f be a numerical function ≥0 defined on T × T', bounded on every compact subset of T × T', and such that: \(1^{\circ}\) for almost every \(t \in \mathrm{T}\) , the function \(t' \mapsto f(t, t')\) is \(\mu'\) -measurable; \(2^{\circ}\) for every function \(h \in \mathcal{K}(\mathrm{T}')\) , the function

\[
t \mapsto \int f (t, t ^ {\prime}) h (t ^ {\prime}) d \mu^ {\prime} (t ^ {\prime}),
\]

defined almost everywhere, is \(\mu\) -measurable. Show that, under these conditions, there exists a function \(g\) , measurable for \(\mu \otimes \mu'\) , such that for every \(t \in \mathrm{T}\) , one has \(f(t, t') = g(t, t')\) except at the points of a \(\mu'\) -negligible set \(\mathrm{A}_t\) (depending on \(t\) ). (Show that for every function \(\varphi \in \mathcal{K}(\mathrm{T} \times \mathrm{T}')\) , the function \(t' \mapsto f(t, t')\varphi(t, t')\) is \(\mu'\) -integrable for almost every \(t \in \mathrm{T}\) , and that the function \(t \mapsto \int f(t, t')\varphi(t, t') d\mu'(t')\) , defined almost everywhere, is \(\mu\) -integrable; one makes use of Lemma 1 of Ch. III, §4, No.~1. Observe next that \(\varphi \mapsto \int d\mu(t) \int f(t, t')\varphi(t, t') d\mu'(t')\) is a positive measure on \(\mathrm{T} \times \mathrm{T}'\) , with base \(\mu \otimes \mu'\) , and apply the Lebesgue-Nikodym theorem; finally, use the fact that there exists a countable dense set in \(\mathcal{K}(\mathrm{T}')\) .)

\begin{enumerate}
\def\labelenumi{\alph{enumi})}
\setcounter{enumi}{1}
\item
  Suppose, moreover, that T is metrizable. Show that the conditions of a) are satisfied if: \(1^{\circ}\) for almost every \(t' \in T'\) , the function \(t \mapsto f(t, t')\) is \(\mu\) -measurable; \(2^{\circ}\) for almost every \(t \in T\) , the function \(t' \mapsto f(t, t')\) is continuous almost everywhere (for \(\mu'\) ). (Use Exer. 13 of Ch. IV, §5.)
\item
  Take \(\mathrm{T} = \mathrm{T}' = [0,1]\) . Admitting the continuum hypothesis (S, Ch. III, §6, No.~4), let \(x \prec y\) be a well-ordering relation on \(\mathrm{T}\) for which there is no greatest element, and such that for every \(x \in \mathrm{T}\) , the set of \(z \prec x\) is countable. If one takes for \(\mu = \mu'\) Lebesgue measure, show that the characteristic function \(f\) of the set of pairs \((t, t')\) such that \(t \prec t'\) satisfies the conditions of \(a\) ), but that \(\int f(t, t') d\mu(t) = 0\) for every \(t' \in \mathrm{T}'\) , and \(\int f(t, t') d\mu'(t') = 1\) for every \(t \in \mathrm{T}\) .
\end{enumerate}

\begin{enumerate}
\def\labelenumi{\arabic{enumi})}
\setcounter{enumi}{7}
\tightlist
\item
  Let \(\mathrm{T}\) be a locally compact space, \(\mu\) a positive measure \(\neq 0\) on \(\mathrm{T}\) , A a discrete space, \(\lambda\) the measure on \(\mathrm{A}\) defined by mass \(+1\) at each point of \(\mathrm{A}\) .
\end{enumerate}

\begin{enumerate}
\def\labelenumi{\alph{enumi})}
\item
  For a mapping \(f\) of \(\mathrm{A} \times \mathrm{T}\) into a topological space to be \((\lambda \otimes \mu)\) -measurable, it is necessary and sufficient that, for every \(\alpha \in \mathrm{A}\) , the mapping \(t \mapsto f(\alpha, t)\) be \(\mu\) -measurable.
\item
  For a function \(\mathbf{f}\) defined on \(\mathrm{A} \times \mathrm{T}\) , with values in \(\overline{\mathbf{R}}\) or in a Banach space, to be \((\lambda \otimes \mu)\) -integrable, it is necessary and sufficient that: \(1^{\circ}\) except for a countable set of values of \(\alpha \in \mathrm{A}\) , the function \(t \mapsto \mathbf{f}(\alpha, t)\) be identically zero; \(2^{\circ}\) for every \(\alpha \in \mathrm{A}\) , the function \(t \mapsto \mathbf{f}(\alpha, t)\) be \(\mu\) -integrable; \(3^{\circ} \sum_{\alpha \in \mathrm{A}} \int |\mathbf{f}(\alpha, t)| d\mu(t) < +\infty\) . For \(\mathbf{f}\) to be
\end{enumerate}

essentially \((\lambda \otimes \mu)\) -integrable, it necessary and sufficient that: \(1^{\circ}\) for each \(\alpha \in \mathbf{A}\) , the function \(t \mapsto \mathbf{f}(\alpha, t)\) be essentially \(\mu\) -integrable; \(2^{\circ} \sum_{\alpha \in \mathbf{A}} \int |\mathbf{f}(\alpha, t)| d\mu(t) < +\infty\) ; one then

has

\[
\iint \mathbf {f} d \lambda d \mu = \sum_ {\alpha \in A} \int \mathbf {f} (\alpha , t) d \mu (t).
\]

\begin{enumerate}
\def\labelenumi{\alph{enumi})}
\setcounter{enumi}{2}
\tightlist
\item
  Let X be a locally compact space, \((\alpha,t)\mapsto\rho_{\alpha,t}\) a family of positive measures on X \((\alpha\in\mathrm{A},t\in\mathrm{T})\) . In order that the family \((\alpha,t)\mapsto\rho_{\alpha,t}\) be \((\mu\otimes\nu)\) -adequate, it is necessary and sufficient that, for every \(\alpha\in\mathrm{A}\) , the family \(t\mapsto\rho_{\alpha,t}\) be \(\mu\) -adequate, and that the family of measures \(\left(\int\rho_{\alpha,t}d\mu(t)\right)_{\alpha\in\mathrm{A}}\) be summable.
\end{enumerate}

\begin{enumerate}
\def\labelenumi{\arabic{enumi})}
\setcounter{enumi}{8}
\tightlist
\item
  Let \(u\) and \(v\) be two increasing and right-continuous numerical functions on \(\mathbf{R}\) , such that \(u(x) = v(x) = 0\) for \(x < 0\) . Let \(w\) be the increasing right-continuous function on \(\mathbf{R}\) defined by \(w(t) = u(t)v(t)\) for \(t \geqslant 0\) , \(w(t) = 0\) for \(t < 0\) ; let \(\lambda, \mu, \nu\) be the Stieltjes measures associated with \(u, v, w\) respectively (§6, Exer. 5).
\end{enumerate}

To every function \(\mathbf{f}\) defined on \(\mathbf{R}\) , with values in \(\overline{\mathbf{R}}\) or in a Banach space \(\mathrm{F}\) , one makes correspond the function \(\overline{\mathbf{f}}\) defined on \(\mathbf{R}^2\) by the conditions \(\overline{\mathbf{f}}(x,y) = \mathbf{f}(x)\) if \(y < x\) , \(\overline{\mathbf{f}}(x,y) = \mathbf{f}(y)\) if \(y \geqslant x\) . Show that, for \(\mathbf{f}\) to be \(\nu\) -integrable, it is necessary and sufficient that \(\overline{\mathbf{f}}\) be integrable for the product measure \(\lambda \otimes \mu\) , in which case \(\int \mathbf{f} d\nu = \iint \overline{\mathbf{f}} d\lambda d\mu\) (prove it first for the characteristic functions of intervals). From this, deduce the formula

\[
\int \mathbf {f} (x) d w (x) = \int \mathbf {f} (x) v (x -) d u (x) + \int \mathbf {f} (x) u (x +) d v (x).
\]

In particular, if \(u\) and \(v\) are continuous on \(\mathbf{R}\) , one obtains the formula for integration by parts

\[
\int_ {a} ^ {b} u (x) d v (x) = u (b) v (b) - u (a) v (a) - \int_ {a} ^ {b} v (x) d u (x).
\]

\begin{enumerate}
\def\labelenumi{\arabic{enumi})}
\setcounter{enumi}{9}
\tightlist
\item
  Let \(\mathbf{T}\) and \(\mathbf{X}\) be two locally compact spaces, \(\lambda\) a positive measure on \(\mathbf{T}\) , \(\mu\) a bounded positive measure on \(\mathbf{X}\) such that \(\mu(\mathbf{X}) = 1\) . Let \(f\) be a numerical function \(> 0\) at every point of \(\mathbf{T} \times \mathbf{X}\) , integrable along with \(\log f\) for the measure \(\lambda \otimes \mu\) . Prove the inequality
\end{enumerate}

\[
\log \left(\int \exp (\int \log f d \mu) d \lambda\right) \leqslant \int (\log \int f d \lambda) d \mu
\]

and show that equality cannot hold unless \(f\) is equivalent to a function of the form \(g \otimes h\) (apply the inequality of the geometric mean (Ch. IV, §6, Exer. 7 d)), for every \(t \in \mathrm{T}\) , to the function \(x \mapsto f(t,x) / (\int f(t,x)d\lambda (t))\) .

\begin{enumerate}
\def\labelenumi{\arabic{enumi})}
\setcounter{enumi}{10}
\tightlist
\item
  Let \(p\) be a finite real number \(\geqslant 1\) , \(T\) and \(X\) two locally compact spaces, \(\lambda\) a positive measure on \(T\) , \(\mu\) a positive measure on \(X\) , and \(f\) a function \(\geqslant 0\) defined on \(T \times X\) , integrable along with \(f^p\) for the measure \(\lambda \otimes \mu\) . Prove the inequality
\end{enumerate}

\[
\left(\int^ {*} \left(\int f d \mu\right) ^ {p} d \lambda\right) ^ {1 / p} \leqslant \int^ {*} \left(\int f ^ {p} d \lambda\right) ^ {1 / p} d \mu .
\]

(For every \(t \in \mathbf{T}\) , apply Hölder's inequality to the function \(x \mapsto f(t, x)\) , put into the form

\[
f (t, x) = g (t, x) \left(\int f ^ {p} (t, x) d \lambda (t)\right) ^ {1 / p q},
\]

where \(q\) is the exponent conjugate to \(p\) .) Show that equality cannot hold unless \(f\) is equivalent to a function of the form \(g \otimes h\) .

¶ 12) Let \(\mathrm{T}_i\) ( \(1 \leqslant i \leqslant n\) ) be \(n\) locally compact spaces, \(\mu_i\) a positive measure on \(\mathrm{T}_i\) for \(1 \leqslant i \leqslant n\) . For each index \(i\) , denote by \(\mathrm{E}_i\) the product \(\prod_{j \neq i} \mathrm{T}_j\) ; let \(f_i\)

be a function \(\geqslant 0\) on \(\mathrm{T} = \prod_{i=1}^{n} \mathrm{T}_i\) , measurable for \(\mu = \bigotimes_{i=1}^{n} \mu_i\) and not dependent

on \(t_{i}\) ; show that if, for \(1 \leqslant k \leqslant n\) , the function \(f_{k}^{n-1}\) is integrable for the measure \(\mu_{1} \otimes \cdots \otimes \mu_{k-1} \otimes \mu_{k+1} \otimes \cdots \otimes \mu_{n}\) , then \(f_{1}f_{2} \cdots f_{n}\) is \(\mu\) -integrable, and

\[
\int f _ {1} f _ {2} \dots f _ {n} d \mu_ {1} d \mu_ {2} \dots d \mu_ {n} \leqslant \left(\prod_ {k = 1} ^ {n} \mathrm{J} _ {k}\right) ^ {1 / (n - 1)}
\]

where, for each index k, one sets \(J_{k}=\int f_{k}^{n-1}d\mu_{1}\cdots d\mu_{k-1}d\mu_{k+1}\cdots d\mu_{n}\) (proceed by induction on n, applying the Lebesgue--Fubini theorem and Hölder's inequality).

Deduce from this that if A is a \(\mu\) -measurable subset of T, \(A_{i}\) its projection on \(E_{i}\) , and if \(A_{i}\) is integrable and of measure \(m_{i}\) (for the measure \(\bigotimes_{j\neq i}\mu_{j}\) on \(E_{i}\) ), then A is

\(\mu\) -integrable and

\[
\mu (\mathrm{A}) \leqslant (m _ {1} m _ {2} \dots m _ {n}) ^ {1 / (n - 1)}.
\]

Examine the case of equality in these two inequalities.

Generalize to the case where, instead of considering the \(n\) products of \(n - 1\) of the \(\mathrm{T}_i\) , one considers the \(\binom{n}{p}\) products of \(p\) of the \(\mathrm{T}_i\) , and where one integrates over \(\mathrm{T}\) a product of \(\binom{n}{p}\) functions \(\geqslant 0\) each of which depends on only \(p\) of the variables \(t_i\) . For example, if \(\mathrm{F}_{ij} = \mathrm{T}_i \times \mathrm{T}_j (i < j)\) and if \(f_{ij}\) depends only on the variables \(t_i\) and \(t_j\) , show that

\[
\int \left(\prod_ {i <   j} f _ {i j}\right) d \mu_ {1} d \mu_ {2} \dots d \mu_ {n} \leqslant \left(\prod_ {i <   j} \int f _ {i j} ^ {n - 1} d \mu_ {i} d \mu_ {j}\right) ^ {1 / (n - 1)}.
\]

¶ 13) Let \((T_{\iota})_{\iota \in I}\) be a family of compact spaces, and for each \(\iota \in I\) , let \(\mu_{\iota}\) be a positive measure on \(T_{\iota}\) of total mass equal to 1. Let \(T\) be the product space \(\prod_{\iota \in I} T_{\iota}\) ,

and \(\mu\) the product measure \(\bigotimes \mu_{\iota}\) on T (Ch. III, §4, No.~5).

\begin{enumerate}
\def\labelenumi{\alph{enumi})}
\item
  For every \(\iota \in I\) , let \(K_{\iota}\) be a compact subset of \(T_{\iota}\) ; show that if \(K = \prod_{\iota \in I} K_{\iota}\) then \(\mu(K) = \prod \mu_{\iota}(K_{\iota})\) .
\item
  Let M be a compact subset of T. Show that there exists a countable subset J of I such that, setting H = I - J, \(T_{J} = \prod_{\iota \in J} T_{\iota}\) , \(\mu_{J} = \bigotimes_{\iota \in J} \mu_{\iota}\) , \(T_{H} = \prod_{\iota \in H} T_{\iota}\) , one has \(M \subset N \times T_{H}\) , where N is a \(\mu_{J}\) -measurable subset of \(T_{J}\) such that \(\mu_{\mathrm{J}}(\overline{\mathrm{N}}) = \mu(\mathrm{M})\) . (Observe that for every n \textgreater{} 0, there exists an open neighborhood of M, the union of a finite number of elementary sets, whose measure differs from \(\mu(\mathrm{M})\) by less than 1/n.) From this, deduce that for almost every \(x \in N\) , the section \(\mathrm{M}(x) \subset \mathrm{T}_{\mathrm{H}}\) (when T is identified with the product \(T_{J} \times T_{H}\) ) contains the product of the supports of the measures \(\mu_{\iota}\) for \(\iota \in H\) (make use of the Lebesgue--Fubini theorem).
\item
  Show that if I is countable and if, for every \(\iota \in I\) , \(A_{\iota}\) is a \(\mu_{\iota}\) -measurable subset of \(T_{\iota}\) , then the set \(A = \prod_{\iota \in I} A_{\iota}\) is \(\mu\) -measurable and \(\mu(A) = \prod_{\iota \in I} \mu_{\iota}(A_{\iota})\) .
\item
  Suppose that I is uncountable. For every \(\iota \in I\) , let \(A_{\iota}\) be a \(\mu_{\iota}\) -measurable subset of \(T_{\iota}\) . For \(A = \prod_{\iota \in I} A_{\iota}\) to be \(\mu\) -measurable, it is necessary and sufficient that one
\end{enumerate}

be in one of the following two cases: \(1^{\circ}\prod_{\iota \in I}\mu_{\iota}(A_{\iota}) = 0;2^{\circ}A_{\iota}\) contains the support of the

measure \(\mu_{\iota}\) , except perhaps for the indices \(\iota\) of a countable subset of I. In each of these two cases, \(\mu(A) = \prod_{\iota \in I} \mu_{\iota}(A_{\iota})\) . (Assuming that one is not in either of the two preceding

cases, show that one can reduce to the case that \(\mu_{\iota}(\mathrm{A}_{\iota}) = 1\) for all \(\iota \in \mathrm{I}\) . Using \(b\) , show that neither A nor T-A can contain a compact set of nonzero measure.)

¶ 14) Let K be the interval {[}0,1{]} of R, λ the Lebesgue measure on K, A an uncountable set, and μ the measure on T = K\^{}A that is the product of the measure λ on each of the factors.

\begin{enumerate}
\def\labelenumi{\alph{enumi})}
\item
  Show that every closed and metrizable subspace X of T is \(\mu\) -negligible (make use of Exer. 13 b) above, and Exer. 8 of GT, IX, §2).
\item
  Let Y be a locally compact space admitting a countable base, \(\nu\) a positive measure on Y, and \(\pi\) a \(\nu\) -proper mapping of Y into T. Show that the image \(\pi(\nu)\) is alien to \(\mu\) (using a), show that \(\pi(\nu)\) is concentrated on a \(\mu\) -negligible set).
\end{enumerate}

\begin{enumerate}
\def\labelenumi{\arabic{enumi})}
\setcounter{enumi}{14}
\tightlist
\item
  Let \((\mathrm{T}_{\iota})_{\iota \in I}\) be any family of compact spaces and, for each \(\iota \in I\) , let \(\mu_{\iota}\) be a positive measure on \(\mathrm{T}_{\iota}\) of total mass 1; let \(\mu\) be the product measure \(\bigotimes_{\iota \in I} \mu_{\iota}\) on \(\mathrm{T} = \prod_{\iota \in \mathrm{I}}\mathrm{T}_{\iota}\) . For every partition (L, M) of I into two subsets, one identifies T with the
\end{enumerate}

product \(\mathrm{T_L} \times \mathrm{T_M}\) , where \(\mathrm{T_L} = \prod_{\iota \in \mathrm{L}} \mathrm{T}_{\iota}\) ; for every \(t \in \mathrm{T}\) , one writes \(t_{\mathrm{L}} = \mathrm{pr}_{\mathrm{L}} t\) , so that

\(t\) is identified with \((t_{\mathrm{L}}, t_{\mathrm{M}})\) ; let \(\mu_{\mathrm{L}}\) be the measure \(\bigotimes_{\iota \in \mathrm{L}} \mu_{\iota}\) on \(\mathrm{T}_{\mathrm{L}}\) . Let \(p\) be a finite

real number \(\geqslant 1\) ; for every function \(\mathbf{f}\) in \(\mathcal{L}_{\mathrm{F}}^{p}(\mathrm{T},\mu)\) (F a Banach space or \(\mathrm{F} = \overline{\mathbf{R}}\) ), one writes \(\mathbf{f}_{\mathrm{L}}(t) = \int \mathbf{f}(t_{\mathrm{L}},t_{\mathrm{M}}) d\mu_{\mathrm{M}}(t_{\mathrm{M}})\) . Show that, with respect to the directed set of finite subsets L of I, \(\mathbf{f}_{\mathrm{L}}\) tends in mean of order \(p\) to \(\mathbf{f}\) , and \(\mathbf{f}_{\mathrm{M}}\) tends in mean of order \(p\) to the constant function equal to \(\int \mathbf{f} d\mu\) . (Approximate \(\mathbf{f}\) by a continuous function depending only on a finite number of variables.)

From this, deduce that if a \(\mu\) -measurable subset \(A\) of \(T\) is such that, for all \(t \in A\) , every point \(t' \in T\) , whose coordinates are equal to those of \(t\) except for a finite number of indices, also belongs to \(A\) , then \(\mu(A) = 0\) or \(\mu(A) = 1\) .

¶ 16) Let \((T_n)\) be an infinite sequence of compact spaces, \(\mu_n\) a positive measure on \(T_n\) of total mass equal to 1, and \(\mu\) the product measure \(\bigotimes_{n=1}^{\infty} \mu_n\) on \(T = \prod_{n=1}^{\infty} T_n\) .

\begin{enumerate}
\def\labelenumi{\alph{enumi})}
\item
  Let \(f\) be a function \(\geqslant 0\) , \(\mu\) -integrable on \(\mathbf{T}\) . Let \((\mathrm{L}_n)\) be an increasing sequence of subsets of \(\mathbf{N}\) , and, setting \(\mathrm{M}_n = \mathbf{N} - \mathrm{L}_n\) , let \(g = \sup_n f_{\mathrm{L}_n}\) , \(h = \sup_n f_{\mathrm{M}_n}\) (the notations of Exer. 15). For every \(\alpha > 0\) , let \(\mathrm{A}_{\alpha}\) be the set of points \(t \in \mathbb{T}\) where \(g(t) > \alpha\) , and \(\mathrm{B}_{\alpha}\) the set of points \(t \in \mathbb{T}\) where \(h(t) > \alpha\) . Show that \(\alpha \cdot \mu(\mathrm{A}_{\alpha}) \leqslant \int f d\mu\) and \(\alpha \cdot \mu(\mathrm{B}_{\alpha}) \leqslant \int f d\mu\) . (Observe that \(\mathrm{A}_{\alpha}\) is the set of \(t \in \mathbb{T}\) where at least one of the \(f_{\mathrm{L}_n}(t)\) is \(> \alpha\) , and express it as a countable union of sets \(\mathrm{G}_n\) , pairwise disjoint and such that \(\alpha \cdot \mu(\mathrm{G}_n) \leqslant \int_{\mathrm{G}_n} f d\mu\) .)
\item
  Suppose that \((\mathrm{L}_n)\) is an increasing sequence of finite subsets of \(\mathbf{N}\) , whose union is \(\mathbf{N}\) . Show that \(f_{\mathrm{L}_n}\) tends almost everywhere to \(f\) , and that \(f_{\mathrm{M}_n}\) tends almost everywhere in \(\mathrm{T}\) to the constant \(\int f d\mu\) . (For every \(\varepsilon > 0\) , consider a continuous function \(g\) depending on only a finite number of variables and such that \(\int |f - g| d\mu \leqslant \varepsilon\) , and apply \(a\) ) to the function \(|f - g|\) .)
\end{enumerate}

\begin{enumerate}
\def\labelenumi{\arabic{enumi})}
\setcounter{enumi}{16}
\tightlist
\item
  Let \((\mathrm{T}_{\iota})_{\iota \in \mathrm{I}}\) be any family of compact spaces; for each \(\iota \in \mathrm{I}\) , let \(\mu_{\iota}\) be a positive measure on \(\mathrm{T}_{\iota}\) , of total mass equal to 1, and let \(\mu\) be the product measure \(\bigotimes_{\iota \in \mathrm{I}} \mu_{\iota}\) on \(\mathrm{T} = \prod_{\iota \in \mathrm{I}} \mathrm{T}_{\iota}\) . For every \(\iota \in \mathrm{I}\) , let \(f_{\iota}\) be a function \(\geqslant 0\) defined on \(\mathrm{T}_{\iota}\) , \(\mu_{\iota}\) -integrable and such that \(\int f_{\iota} d\mu_{\iota} \leqslant 1\) ; set \(\mu_{\iota}' = f_{\iota} \cdot \mu_{\iota}\) and \(\mu' = \bigotimes_{\iota \in \mathrm{I}} \mu_{\iota}'\) (Ch. III, §4, Exer. 6). Assume that \(\mu' \neq 0\) .
\end{enumerate}

\begin{enumerate}
\def\labelenumi{\alph{enumi})}
\tightlist
\item
  For every finite subset \(J\) of \(I\) , set \(T_J = \prod_{\iota \in J} T_\iota\) , \(\mu_J = \bigotimes_{\iota \in J} \mu_\iota\) , \(\mu_J' = \bigotimes_{\iota \in J} \mu_\iota'\) ,
\end{enumerate}

\(f_{\mathrm{J}}(t) = \prod_{\iota \in \mathbb{J}}f_{\iota}(\mathrm{pr}_{\iota}t), g_{\mathrm{J}} = \sqrt{f_{\mathrm{J}}}, \mu_{\mathrm{J}}^{\prime \prime} = \sqrt{\mu_{\mathrm{J}}\mu_{\mathrm{J}}^{\prime}}\) (§5, No.~9) and \(\rho (\mu_{\mathrm{J}},\mu_{\mathrm{J}}^{\prime}) = \mu_{\mathrm{J}}^{\prime \prime}(\mathrm{T}_{\mathrm{J}}) =\)

\(\int g_{\mathrm{J}}d\mu_{\mathrm{J}}\) . Show that, for the functions \(g_{\mathrm{J}}\) to converge in the quadratic mean to a function in \(\mathcal{L}^2 (\mathrm{T},\mu)\) , with respect to the directed ordered set of finite subsets of I, it is necessary and sufficient that the product of the numbers \(\rho (\mu_{\iota},\mu_{\iota}^{\prime})\) be convergent in \(\mathbf{R}_{+}^{*}\) , that is, \(>0\) . (For two finite subsets J,L such that \(\mathrm{J}\subset \mathrm{L}\) , evaluate the norm \(\mathrm{N}_2(g_{\mathrm{J}} - g_{\mathrm{L}})\) by means of the \(\rho (\mu_{\iota},\mu_{\iota}^{\prime})\) .) Deduce from this that the functions \(f_{\mathrm{J}}\) then converge in mean to a function \(f\geqslant 0\) such that \(\mu^{\prime} = f\cdot \mu\) .

\begin{enumerate}
\def\labelenumi{\alph{enumi})}
\setcounter{enumi}{1}
\tightlist
\item
  Show that if the product \(\prod_{\iota \in I} \rho(\mu_{\iota}, \mu_{\iota}')\) is zero, then the measures \(\mu\) and \(\mu'\) are
\end{enumerate}

alien (consider the set \(\mathrm{A_J}\) of points \(t\in \mathrm{T}\) such that \(g_{\mathrm{J}}(t) > 1\) , and evaluate the measures \(\mu (\mathrm{A_J})\) and \(\mu '(T - A_J)\) ).

¶ 18) The notations being the same as in Exer. 17, let \(\nu_{\iota}\) be a positive measure on \(\mathrm{T}_{\iota}\) , of total mass equal to 1, and let \(\nu = \bigotimes \nu_{\iota}\) .

\begin{enumerate}
\def\labelenumi{\alph{enumi})}
\item
  Show that if one of the measures \(\nu_{\iota}\) is alien to the measure \(\mu_{\iota}\) with the same index, then \(\nu\) is alien to \(\mu\) .
\item
  For every \(\iota \in I\) , write \(\nu_{\iota} = \mu_{\iota}' + \mu_{\iota}''\) , where \(\mu_{\iota}'\) has base \(\mu_{\iota}\) , and \(\mu_{\iota}''\) is alien to \(\mu_{\iota}\) (§5, Th. 3); suppose that \(\mu_{\iota}' \neq 0\) for all \(\iota \in I\) . Show that, for \(\nu\) not to be alien to \(\mu\) , it is necessary and sufficient that \(\mu_{\iota}'' = 0\) except for a countable family of indices, and that the product \(\mu' = \bigotimes_{\iota \in I} \mu_{\iota}'\) be a measure \(\neq 0\) with base \(\mu\) ; \(\mu'' = \nu - \mu'\) is then
\end{enumerate}

alien to \(\mu\) . (If \(\mu_{\iota}^{\prime\prime} \neq 0\) for an uncountable infinity of indices, show that there exists a number \(\alpha\) such that \(0 < \alpha < 1\) and such that, for a countable infinity of indices \(\iota_{n}\) , one has \(\mu_{\iota_{n}}^{\prime}(T_{\iota_{n}}) \leqslant \alpha\) ; show then that \(\mu\) is alien to \(\nu\) by using \(a\) ); prove the second part of the proposition by using Exer. 17.)

\begin{enumerate}
\def\labelenumi{\arabic{enumi})}
\setcounter{enumi}{18}
\tightlist
\item
  Let \(X\) be a locally compact space, \(Y\) a locally compact space countable at infinity, \(A\) a universally measurable subset of \(X \times Y\) .
\end{enumerate}

\begin{enumerate}
\def\labelenumi{\alph{enumi})}
\item
  For every \(x \in X\) , show that the section \(\mathrm{A}(x)\) is a universally measurable subset of Y. Moreover, for every positive measure \(\mu\) on Y, the function \(x \mapsto \mu^{*}(\mathrm{A}(x))\) is universally measurable on X (make use of the Lebesgue--Fubini theorem).
\item
  Let \(\mu\) be a positive measure on Y such that, for almost every \(y \in Y\) (for \(\mu\) ), the section \(\overline{A}(y)\) of A is countable. Show that the set N of \(x \in X\) such that \(\mu^{*}(A(x)) > 0\) cannot contain any countable compact set.
\end{enumerate}

\begin{enumerate}
\def\labelenumi{\arabic{enumi})}
\setcounter{enumi}{19}
\tightlist
\item
  \begin{enumerate}
  \def\labelenumii{\alph{enumii})}
  \tightlist
  \item
    Let \(\nu\) be a positive measure on the unit circle \(\mathbf{U}: |z| = 1\) in \(\mathbf{C}\) , and let \(\mu\) be the complex measure \(j \cdot \nu\) , where \(j: z \mapsto z\) is the canonical injection of \(\mathbf{U}\) into \(\mathbf{C}\) , so that \(|\mu| = \nu\) . Let \(\rho\) be the inverse image of \(\nu\) under the canonical local homeomorphism \(\varphi: t \mapsto e^{it}\) of \(\mathbf{R}\) onto \(\mathbf{U}\) (§6, No.~6). Let A be a \(\nu\) -measurable subset of \(\mathbf{U}\) ; if \(\mu(A) = |\mu(A)|e^{i\omega}\) , show that \(|\mu(A)| \leqslant \left|\int_{\omega - \pi/2}^{\omega + \pi/2} e^{i\theta} d\rho(\theta)\right|\) .
  \end{enumerate}
\end{enumerate}

\begin{enumerate}
\def\labelenumi{\alph{enumi})}
\setcounter{enumi}{1}
\tightlist
\item
  Deduce from \(a\) ) that, as A runs over the set of \(\mu\) -measurable subsets of \(\mathbf{U}\) , one has
\end{enumerate}

\[
\begin{array}{l} \sup _ {\mathrm{A}} | \mu (\mathrm{A}) | \geqslant \sup _ {\omega} \left| \int_ {\omega - \pi / 2} ^ {\omega + \pi / 2} \mathcal {R} \big (e ^ {i (\theta - \omega)} \big) d \rho (\theta) \right| \\ \geqslant \frac {1}{2 \pi} \int_ {0} ^ {2 \pi} d \omega \int_ {\omega - \pi / 2} ^ {\omega + \pi / 2} \mathcal {R} \big (e ^ {i (\theta - \omega)} \big) d \rho (\theta) \end{array}
\]

and conclude therefrom that

\[
\sup _ {A} | \mu (A) | \geqslant \frac {1}{\pi} \| \mu \|\tag{*}
\]

(evaluate the last integral with the help of the Lebesgue--Fubini theorem).

\begin{enumerate}
\def\labelenumi{\alph{enumi})}
\setcounter{enumi}{2}
\tightlist
\item
  Show that the inequality (*) is true for every bounded complex measure on a locally compact space X, where A runs over the set of \(\mu\) -measurable subsets of X. (Let \(\mu = h \cdot |\mu|\) , where h takes its values in U. Consider the image measure \(\nu = h(|\mu|)\) on U and show that \(h(\mu) = j \cdot \nu\) ; then apply b) to \(h(\mu)\) .)
\end{enumerate}

\begin{enumerate}
\def\labelenumi{\arabic{enumi})}
\setcounter{enumi}{20}
\tightlist
\item
  Denote by \(\lambda\) the Lebesgue measure on \(\mathbf{R}^n\) .
\end{enumerate}

\begin{enumerate}
\def\labelenumi{\alph{enumi})}
\tightlist
\item
  Let \(\mathbf{C}_0\) be a cube (GT, VI, §1, No.~1) in \(\mathbf{R}^n\) , and let \(u\) be a numerical function on \(\mathbf{C}_0\) that is integrable (for \(\lambda\) ). Let \(s\) be a real number such that
\end{enumerate}

\[
s \geqslant \frac {1}{\lambda (\mathrm{C} _ {0})} \int_ {\mathrm{C} _ {0}} | u | d \lambda .
\]

Show that there exists a sequence \((\mathrm{I}_k)\) of open cubes contained in \(\mathbf{C}_0\) , pairwise disjoint, and satisfying the following conditions:

\(1^{\circ}\left|u(x)\right|\leqslant s\) almost everywhere in the set \(\mathrm{C_0 - }\bigcup_{k}\mathrm{I}_k;\)

\(2^{\circ}\) if one sets

\[
u _ {k} = \frac {1}{\lambda (\mathrm{I} _ {k})} \int_ {\mathrm{I} _ {k}} | u (x) | d \lambda (x),
\]

then \(u_{k} \leqslant 2^{n}s\) for all \(k\) ;

\[
3 ^ {\circ} \sum_ {k} \lambda (I _ {k}) \leqslant \frac {1}{s} \int_ {C _ {0}} | u (x) | d \lambda (x).
\]

(Dividing into two equal intervals each of the projections of \(C_{0}\) , among the \(2^{n}\) pairwise disjoint open cubes choose the cubes \(I_{1i}\) such that the corresponding `mean values' \(u_{1i}\) are all \(\geqslant s\) , and, because of the choice of s, show that \(u_{1i} \leqslant 2^{n}s\) . Similarly decompose into \(2^{n}\) cubes each of the cubes of the first subdivision that does not figure among the \(I_{1i}\) , which yields a second finite family of cubes \(I_{2i}\) chosen among all of the new cubes obtained as those for which the corresponding mean value \(u_{2i}\) is \(\leqslant s\) . Continue by recursion. To prove that the condition \(1^{\circ}\) is verified, argue by contradiction, noting that the subset \(B_{m}\) of \(C_{0} - \bigcup_{k} I_{k}\) where \(|u(x)| \geqslant s + 1/m\) is, up to a negligible set, the intersection of a decreasing sequence \((\mathrm{U}_r)\) of open sets, where \(\mathrm{U}_r\) is the union of the cubes of the \(r\) -th subdivision that do not figure among the \(\mathrm{I}_k\) and that contain a point of \(\mathrm{B}_m\) .

\begin{enumerate}
\def\labelenumi{\alph{enumi})}
\setcounter{enumi}{1}
\tightlist
\item
  Let \(\mathcal{F}\) be the set of integrable numerical functions \(u\) on \(C_0\) having the following property: for every cube \(C \subset C_0\) , if \(m_C(u)\) is the mean value
\end{enumerate}

\[
\frac {1}{\lambda (\mathrm{C})} \int_ {\mathrm{C}} u (x) d \lambda (x),
\]

then \(\int_{\mathbb{C}} |u(x) - m_{\mathbb{C}}(u)| d\lambda(x) \leqslant \lambda(\mathbb{C})\) ; if \(u \in \mathcal{F}\) , the same is true of \(u - c\) for every real number \(c\) . For every function \(u \in \mathcal{F}\) , every number \(\sigma > 0\) and every cube \(\mathbb{C} \subset \mathbb{C}_0\) , denote by \(\mathrm{S}(\sigma, u, \mathbb{C})\) the set of \(x \in \mathbb{C}\) such that \(|u(x) - m_{\mathbb{C}}(u)| \geqslant \sigma\) . Finally, denote by \(\mathrm{G}(\sigma)\) the smallest number such that

\[
\lambda \bigl (\mathrm{S} (\sigma , u, \mathrm{C}) \bigr) \leqslant \mathrm{G} (\sigma) \int_ {\mathrm{C}} | u (x) - m _ {\mathrm{C}} (u) | d \lambda (x)
\]

as \(u\) runs over \(\mathcal{F}\) , and \(C\) over the set of cubes contained in \(C_0\) ; one has \(G(\sigma) \leqslant 1 / \sigma\) . Show that if \(\sigma / 2^n > s \geqslant 1\) then

\[
\mathrm{G} (\sigma) \leqslant \frac {1}{s} \mathrm{G} (\sigma - 2 ^ {n} s).
\]

(One can restrict oneself to proving that, for a \(u \in \mathcal{F}\) , one has

\[
\lambda \left(\mathrm{S} (\sigma , u, \mathrm{C} _ {0})\right) \leqslant \frac {1}{s} \mathrm{G} (\sigma - 2 ^ {n} s) \int_ {\mathrm{C} _ {0}} | u (x) - m _ {\mathrm{C} _ {0}} (u) | d \lambda (x),
\]

and one can suppose, to simplify, that \(m_{\mathrm{C}_0}(u) = 0\) . Make use then of the decomposition of \(\mathrm{C}_0\) into a union of cubes \(\mathrm{I}_k\) and a set where \(|u(x)| \leqslant s\) almost everywhere, on noting that if \(x \in \mathrm{I}_k\) belongs to \(\mathrm{S}(\sigma, u, \mathrm{C}_0)\) , then \(|u(x) - m_{\mathrm{I}_k}(u)| \geqslant \sigma - 2^n s\) and consequently

\[
\lambda \left(\mathrm{I} _ {k} \cap \mathrm{S} (\sigma , u, \mathrm{C} _ {0})\right) \leqslant \mathrm{G} (\sigma - 2 ^ {n} s) \int_ {\mathrm{I} _ {k}} | u (x) - m _ {\mathrm{I} _ {k}} (u) | d \lambda (x). \Big)
\]

\begin{enumerate}
\def\labelenumi{\alph{enumi})}
\setcounter{enumi}{2}
\tightlist
\item
  Deduce from \(b\) ) that there exist two constants \(B, b\) depending only on \(n\) , such that, for every function \(u \in \mathcal{F}\) and every number \(\sigma > 0\) , one has \(\lambda(\mathrm{S}(\sigma, u, \mathrm{C}_0)) \leqslant \mathrm{Be}^{-b\sigma}\lambda(\mathrm{C}_0)\) . (Take \(s = e\) in \(b\) ).)
\end{enumerate}

